% Einheit 4 von 4 – Prozentuale Veränderung (Katalog Einheit 5)
\einheitenkopf[e4][4 Veränderung]{Einheit 4 von 4 · Prozentuale Veränderung}
\zweigzeile{Hier lernst du auszurechnen, wie sich ein Wert ändert, wenn er um einige Prozent steigt oder fällt · neu in diesem Jahr oder erst in Klasse 8 – je nach Buch (am Gymnasium in Klasse 7) · P10 oft · baut auf: Prozentsatz (Einheit 1), Prozentwert (Einheit 2), Grundwert (Einheit 3), Runden (Nr.~6)}
\setcounter{aufgabe}{41}

\begin{aufgabe}{Ich erkenne, ob etwas um einen Prozentsatz oder auf einen Prozentsatz geändert wird.}
Kreuze an, was dasselbe bedeutet.
\begin{teile}
\teil Der Preis wird um $20\,\%$ gesenkt.\\ \kreuz{auf $20\,\%$ gesenkt}\quad \kreuz{auf $80\,\%$ gesenkt}\quad \kreuz{auf $120\,\%$ gesenkt}
\teil Die Miete wird um $10\,\%$ erhöht.\\ \kreuz{auf $110\,\%$ erhöht}\quad \kreuz{auf $10\,\%$ erhöht}\quad \kreuz{auf $90\,\%$ erhöht}
\teil Die Jacke ist auf $75\,\%$ herabgesetzt.\\ \kreuz{um $75\,\%$ herabgesetzt}\quad \kreuz{um $25\,\%$ herabgesetzt}
\teil Der Eintritt steigt auf $125\,\%$.\\ \kreuz{um $125\,\%$ erhöht}\quad \kreuz{um $25\,\%$ erhöht}
\teil Der Preis wird auf $60\,\%$ gesenkt.\\ \kreuz{um $40\,\%$ gesenkt}\quad \kreuz{um $60\,\%$ gesenkt}
\end{teile}
\end{aufgabe}

\begin{aufgabe}{Ich erkenne, welcher Wert der alte ist.}
Unterstreiche den alten Wert.
\begin{teile}
\teil Der Preis stieg von $40\,€$ auf $50\,€$.
\teil Jetzt kostet die Karte $12\,€$. Vorher kostete sie $10\,€$.
\teil Der Zug braucht nur noch 45 Minuten statt 60 Minuten.
\teil 2025 kamen 800 Besucher. Im Jahr davor waren es 1000.
\teil Nach der Erhöhung zahlt Lina $330\,€$ Miete. Bisher waren es $300\,€$.
\end{teile}
\end{aufgabe}

\verfahren{Neuer Wert: Teil dazu oder weg}
\begin{aufgabe}{Ich kann ausrechnen, wie groß ein Wert nach einer Erhöhung oder Senkung ist.}
Berechne den neuen Preis.

\swz{$60\,€$ werden um $10\,\%$ teurer.}{\streifen{10}{$0\,€$}{$60\,€$}}
\swz{$60\,€$ werden um $20\,\%$ teurer.}{\streifen{20}{$0\,€$}{$60\,€$}}
\swz{$60\,€$ werden um $50\,\%$ teurer.}{\streifen{50}{$0\,€$}{$60\,€$}}
\swz{$60\,€$ werden um $25\,\%$ teurer.}{\streifen[4]{25}{$0\,€$}{$60\,€$}}
\swfrage{Was bleibt gleich, was ändert sich?}
\end{aufgabe}

\begin{aufgabe}{Ich kann ausrechnen, wie groß ein Wert nach einer Erhöhung oder Senkung ist – weiter.}
Berechne den neuen Preis.

\swz{$40\,€$ werden um $10\,\%$ billiger.}{\streifen{10}{$0\,€$}{$40\,€$}}
\swz{$90\,€$ werden um $30\,\%$ billiger.}{\streifen{30}{$0\,€$}{$90\,€$}}
\end{aufgabe}

\verfahren{Neuer Wert mit dem Faktor}
\begin{aufgabe}{Ich kann den neuen Wert in einem Schritt ausrechnen (kommt in Klasse 9; am Gymnasium schon in Klasse 7).}
Rechne mit dem Faktor. Diese Zahl nennt man Wachstumsfaktor.

\swz{$40\,€$ um $20\,\%$ erhöht}{$100\,\% + 20\,\% = 120\,\% = 1{,}2$}
\swz{$70\,€$ um $20\,\%$ gesenkt}{$100\,\% - 20\,\% = 80\,\% = 0{,}8$}
\swz{$200\,€$ um $5\,\%$ erhöht}{$100\,\% + 5\,\% = \rule{10mm}{0.4pt}\,\% = \rule{10mm}{0.4pt}$}
\end{aufgabe}

\verfahren{Veränderung in Prozent}
\begin{aufgabe}{Ich kann ausrechnen, um wie viel Prozent sich ein Wert verändert hat.}
Berechne, um wie viel Prozent sich der Wert verändert hat.

\swz[3]{Ein Preis steigt von $50\,€$ auf $60\,€$.}{\streifen{0}{$0\,€$}{$50\,€$}}
\swz[3]{Ein Preis sinkt von $80\,€$ auf $60\,€$.}{\streifen{0}{$0\,€$}{$80\,€$}}
\swz[4]{Ein Liter Milch kostete im Januar $1{,}15\,€$, im Juni $1{,}22\,€$. Um wie viel Prozent ist der Preis gestiegen? Runde auf eine Stelle nach dem Komma. (P10 2026 FOR)}{}
\end{aufgabe}

\verfahren{Alter Wert aus neuem Wert}
\begin{aufgabe}{Ich kann aus dem neuen Wert den alten ausrechnen.}
Berechne den alten Preis.

\swz[2]{Nach $20\,\%$ Rabatt kostet ein Spiel $32\,€$.}{\begin{dreisatz}{Prozent}{Euro}
\dsz{80\,\%}{32\,€}
\dsp{:80}{:80}
\dsz{1\,\%}{\dsleer}
\dsp{\cdot 100}{\cdot 100}
\dsz{100\,\%}{\dsleer}
\end{dreisatz}}
\swz[2]{Nach einer Erhöhung um $25\,\%$ kostet der Eintritt $10\,€$.}{\begin{dreisatz}{Prozent}{Euro}
\dsz{\rule{8mm}{0.4pt}\,\%}{10\,€}
\dsp{\phantom{:0}}{\phantom{:0}}
\dsz{1\,\%}{\dsleer}
\dsp{\phantom{:0}}{\phantom{:0}}
\dsz{100\,\%}{\dsleer}
\end{dreisatz}}
\end{aufgabe}

\verfahren{Brutto und Netto}
\begin{aufgabe}{Ich kann Preise mit und ohne Mehrwertsteuer ausrechnen (neu in diesem Jahr).}
Brutto heißt mit Mehrwertsteuer, netto heißt ohne. Rechne.

\swz{Netto $200\,€$, dazu $19\,\%$. Wie viel ist brutto?}{\streifen{100}{$0\,€$}{$200\,€$}}
\swz{Netto $30\,€$, dazu $7\,\%$. Wie viel ist brutto?}{\streifen{100}{$0\,€$}{$30\,€$}}
\swz[2]{Brutto $59{,}50\,€$ mit $19\,\%$. Wie viel ist netto?}{\begin{dreisatz}{Prozent}{Euro}
\dsz{119\,\%}{59{,}50\,€}
\dsp{\phantom{:0}}{\phantom{:0}}
\dsz{1\,\%}{\dsleer}
\dsp{\phantom{:0}}{\phantom{:0}}
\dsz{100\,\%}{\dsleer}
\end{dreisatz}}
\end{aufgabe}

\verfahren{Prozentpunkte}
\begin{aufgabe}{Ich kann Prozentpunkte und Prozent unterscheiden.}
Der Unterschied zweier Prozentangaben heißt Prozentpunkte. Rechne.

\swz{Die Wahlbeteiligung stieg von $40\,\%$ auf $50\,\%$. Um wie viele Prozentpunkte?}{\streifen{50}{$0\,\%$}{$100\,\%$}}
\swz{Um wie viel Prozent ist die Wahlbeteiligung gestiegen?}{\streifen[4]{0}{$0$}{$40\,\%$ (alt)}}
\swz[3]{In einer Schule mit 600 Kindern stieg der Anteil der Radfahrer von $20\,\%$ auf $25\,\%$. Wie viele Kinder fahren jetzt mehr mit dem Rad?}{\streifen[20]{25}{$0$}{600 Kinder}}
\end{aufgabe}

\verfahren{Steigung in Prozent}
\begin{aufgabe}{Ich kann eine Steigung in Prozent ausrechnen (kommt in Klasse 10).}
Die Steigung in Prozent sagt, wie viele Meter es auf 100 m waagerechter Strecke hochgeht. Rechne.

\swz{Ein Schild zeigt $12\,\%$ Steigung. Wie viele Meter steigt die Straße auf $100\,$m waagerechter Strecke?}{\dreieckrw{6}{1.2}{100\text{ m}}{\text{?}}{}}
\swz{Ein Weg steigt auf $60\,$m waagerechter Strecke um $3\,$m. Wie groß ist die Steigung in Prozent?}{\dreieckrw{6}{1}{60\text{ m}}{3\text{ m}}{}}
\swz[4]{Ein Radweg darf höchstens $4\,\%$ Steigung haben. Ein geplanter Weg steigt auf $280\,$m waagerechter Strecke um $13\,$m. Berechne die Steigung auf eine Stelle nach dem Komma. Ist der Weg erlaubt? (P10 2025 OS)}{}
\end{aufgabe}

\begin{aufgabe}{Ich finde den Fehler, wenn jemand eine Veränderung in Prozent ausrechnet.}
Finde den Fehler, benenne ihn und rechne richtig.

\swz[3]{Der Preis stieg von $20\,€$ auf $25\,€$. Um wie viel Prozent? Nils rechnet so:}{\rechnung{25 - 20 &= 5 \\ 5 : 25 &= 0{,}2 = 20\,\%}}
\swz[3]{Der Preis sank von $50\,€$ auf $40\,€$. Um wie viel Prozent? Aylin rechnet so:}{\rechnung{50 - 40 &= 10 \\ 10 : 40 &= 0{,}25 = 25\,\%}}
\end{aufgabe}

\begin{aufgabe}{Ich kann begründen, welcher Wert 100 \% ist.}
Begründe.

\swfrage{Bei „um $20\,\%$ teurer“: Welcher Preis ist $100\,\%$, der alte oder der neue?}
\swfrage{Ein Preis steigt von $40\,€$ auf $50\,€$ und fällt danach wieder von $50\,€$ auf $40\,€$. Sind beide Veränderungen gleich viel Prozent? Begründe ohne zu rechnen.}
\end{aufgabe}

\begin{aufgabe}{Ich kann zwei Rabattangebote vergleichen.}
Eine Jacke kostet $100\,€$. Laden A senkt den Preis erst um $20\,\%$. Eine Woche später senkt er den neuen Preis noch einmal um $10\,\%$. Laden B senkt den Preis sofort um $30\,\%$. Rechne und entscheide.

\swz[3]{Wie viel kostet die Jacke am Ende in Laden A?}{\streifen{20}{$0\,€$}{$100\,€$}}
\swz{Wie viel kostet die Jacke in Laden B?}{\streifen{30}{$0\,€$}{$100\,€$}}
\swfrage{Welcher Laden ist günstiger? Warum sind $20\,\%$ und dann $10\,\%$ nicht dasselbe wie $30\,\%$?}
\end{aufgabe}

\uebersichtskasten{Erhöhung um p \%: neuer Wert = alter Wert $\cdot\,(1 + p/100)$; Senkung: $\cdot\,(1 - p/100)$. \quad 250 € um 12 \% erhöht: $250 \cdot 1{,}12 = 280$ € \quad 250 € um 12 \% gesenkt: $250 \cdot 0{,}88 = 220$ € \\ Veränderung in Prozent: (neu $-$ alt) : alt. \quad von 60 auf 69: $9 : 60 = 0{,}15 = 15$ \% \\ Prozentpunkte: Unterschied zweier Prozentangaben. \quad von 30 \% auf 33 \%: 3 Prozentpunkte (das sind 10 \% mehr) \\ Steigung in Prozent: Höhe : waagerechte Strecke. \quad 8 m auf 100 m = 8 \%}
